\chapter[Literature Review]{Literature Review}{Literature Review}\label{lit}
\noindent
\rule{0.49\textwidth}{0.75pt} $_{\bigcirc}$ \rule{0.49\textwidth}{0.75pt}\\

\begingroup
\linespread{0.8}\selectfont
\section{FedRFQ: Prototype-Based Federated Learning With Reduced Redundancy, Minimal Failure and Enhanced Quality}
\endgroup
FedRFQ presents a prototype-based federated learning approach that reduces redundancy while improving the robustness and quality of local updates. The method is useful because it shows that representative summaries of client updates can be leveraged instead of exchanging full, repetitive information at every round. However, the framework still requires larger-scale validation under varied non-IID settings and does not fully address secure auditability or incentive-aware defense.

DAVS builds on this insight by using compressed gradient sketches to retain representative update information while lowering communication cost. In addition, DAVS extends the idea of representation-based selection beyond compression by using the current-round sketch statistics to identify trustworthy verifiers and support secure aggregation.

\begingroup
\linespread{0.8}\selectfont
\section{LearningChain: A Highly Scalable and Applicable Learning-Based Blockchain Performance Optimization Framework}
\endgroup
LearningChain introduces a learning-driven framework for optimizing blockchain performance through workload-aware tuning, profiling, and performance prediction. The work is valuable because it highlights that blockchain systems can be improved through adaptive control rather than fixed configuration. Its main limitation is that the integration of such optimization ideas into privacy-sensitive federated learning workflows remains underexplored.

DAVS differs by focusing not on tuning blockchain performance alone, but on combining blockchain with federated learning for secure medical training. The proposed framework uses the blockchain as an immutable audit layer, while the communication and validation logic is handled by gradient sketching, verifier selection, and PBFT-style consensus.

\begingroup
\linespread{0.8}\selectfont
\section{PrestigeBFT: Revolutionizing View Changes in BFT Consensus with Reputation Mechanisms}
\endgroup
PrestigeBFT improves BFT liveness and robustness by using reputation-aware view-change procedures. The approach demonstrates that consensus systems can be made more adaptive by using node behavior to guide leader-related decisions. However, the reliance on reputation creates a new challenge: the reputation itself must be bootstrapped and protected from manipulation in adversarial settings.

DAVS avoids this weakness by remaining amnesic and evaluating only the current round of gradient sketches. Instead of depending on historical behavior, DAVS selects verifiers based on similarity and magnitude consistency, which makes the system more resistant to delayed poisoning and trust-building attacks.

\begingroup
\linespread{0.8}\selectfont
\section{BlockFL: A Blockchain-Enabled Federated Learning System for Securing IoVs}
\endgroup
BlockFL is a blockchain-enabled federated learning system that supports secure aggregation and provenance tracking in vehicular networks. It shows how blockchain can be used to make FL more trustworthy, especially in distributed environments. At the same time, the framework still faces communication and latency constraints in highly mobile settings, and its real-world deployment evaluation remains open.

DAVS applies the blockchain-enabled FL idea to medical imaging instead of vehicles and strengthens the design with communication-efficient sketching and local verifier selection. This allows DAVS to address both security and bandwidth limitations in a more controlled healthcare environment.

\begingroup
\linespread{0.8}\selectfont
\section{BlockDFL: A Blockchain-Based Fully Decentralized Peer-to-Peer Federated Learning Framework}
\endgroup
BlockDFL explores fully decentralized federated learning with blockchain-based coordination. Its contribution lies in showing that peer-to-peer aggregation and smart-contract-based model registry can remove the dependence on a central server. However, the framework still faces recurring challenges in deployment complexity, consensus overhead, and communication efficiency.

DAVS is aligned with the same decentralization goal, but narrows the focus to secure medical FL where committee selection and PBFT consensus are used only after sketch-based reduction of communication overhead. This makes the proposed framework more practical for a bandwidth-sensitive clinical setting.

\begingroup
\linespread{0.8}\selectfont
\section{DeepChain: A Deep Learning + Blockchain Framework for Detecting Risky Transactions on HIE Systems}
\endgroup
DeepChain integrates deep learning with blockchain-based immutable logging to support auditable detection in health-information exchange systems. The framework is useful because it combines machine intelligence with traceable records for sensitive workflows. Its limitation is that combining heavy deep learning inference with blockchain logging raises scalability concerns, and privacy-preserving training is not addressed in depth.

DAVS focuses specifically on privacy-preserving distributed training for medical imaging. It avoids exposing raw data during training by combining local learning, gradient sketching, and blockchain-backed auditability in a way that is more directly suited to federated model development.

\begingroup
\linespread{0.8}\selectfont
\section{Biscotti: A Blockchain System for Private and Secure Federated Learning}
\endgroup
Biscotti is one of the early systems to combine blockchain with federated learning for privacy, secure aggregation, and auditability. It establishes the usefulness of blockchain-based coordination in collaborative learning. However, communication and latency overhead remain major limitations, and the incentive design is not fully resolved.

DAVS shares the same privacy and auditability goals, but adds a stronger round-local defense mechanism by using data-aware verifier selection. The use of gradient sketches also reduces communication cost more aggressively than full-gradient blockchain-based systems.

\begingroup
\linespread{0.8}\selectfont
\section{RepChain: Reputation-Based Secure, Fast and High-Incentive Blockchain via Sharding}
\endgroup
RepChain uses reputation scoring with sharding to improve blockchain throughput and incentivize honest participation. The framework is important because it demonstrates how blockchain performance can be improved with protocol-level design choices. Its remaining challenges include reputation manipulation, observe-act adversaries, and the need for further dynamic parameter tuning.

DAVS adopts the broader security lesson from RepChain but does not depend on reputation. Instead, it uses current-round similarity scoring for committee selection, which makes the validator set more difficult to manipulate in adaptive attack scenarios.

\subsection{Problem Statement}
Despite the progress made by the above systems, existing blockchain and federated learning approaches still face major issues such as poisoning attacks, communication-heavy training, lack of robust client selection, and limited traceability. These shortcomings motivate DAVS, which combines compression, committee validation, consensus, and immutable logging into one secure pipeline.

\subsection{Research Gap}
The literature shows that existing works typically solve only one part of the problem: some reduce redundancy, some improve blockchain performance, some strengthen consensus, and some improve auditability. What is still missing is a single communication-efficient framework that can jointly support robust verifier selection, Byzantine resilience, and audit-ready federated learning for sensitive medical data.

\rule{0.49\textwidth}{0.75pt} $_{\bigcirc}$ \rule{0.49\textwidth}{0.75pt}\\
